% GENERATED FILE. Edit canonical lessons, then run npm run generate:books.
% canonical-source-hash: fnv1a64:fec8111cfea4b7df
% canonical-lessons: KA-C120-galipata, KA-C120-angi, KA-C120-sire, KA-C120-pamce, KA-C120-cappali

\chapter{Kite, Shirt, Sari, Dhoti, Sandals}
\label{ch:ka-kite-shirt-sari-dhoti-sandals}
\hlchaptermodality{120}

\begin{chapteropening}
\textbf{By the end of this chapter:} \emph{I can say a kite, a shirt, a sari, a dhoti and sandals.}

Complete the last lesson of chapter 120: i can say a kite, a shirt, a sari, a dhoti and sandals.
\end{chapteropening}

\section[\texorpdfstring{\kn{ಗಾಳಿಪಟ} (gāḷipaṭa)}{ಗಾಳಿಪಟ (gāḷipaṭa)}]{\kn{ಗಾಳಿಪಟ} (gāḷipaṭa) — a kite}

\label{lesson:KA-C120-galipata}



\begin{quote}
Before the new one: say the Kannada for a doll, then the Kannada for a ball.
\end{quote}

\subsection*{You'll want to know: \kn{ಗಾಳಿಪಟ}}
\textbf{\kn{ಗಾಳಿಪಟ}} — \emph{gāḷipaṭa} — \textquotedblleft{}a kite\textquotedblright{}.

\begin{grammarlens}[title={What you've built}]
One more word for this chapter.
\end{grammarlens}

\subsection*{Your turn}
\begin{itemize}
\raggedright
  \item \emph{Say it:} \emph{gāḷipaṭa}
  \item \emph{Say it:} \emph{gāḷipaṭa}, once more
  \item \emph{Say it:} say \emph{gāḷipaṭa}
  \item \emph{Recall:} say the Kannada for a doll, then the Kannada for a ball, then say \emph{gāḷipaṭa} again
\end{itemize}

\begin{tcolorbox}[breakable,colback=teal!4,colframe=teal!35!black,title={Before you move on}]
What does \kn{ಗಾಳಿಪಟ} mean? (\textquotedblleft{}A kite\textquotedblright{}.) Say it once more.
\end{tcolorbox}
\section[\texorpdfstring{\kn{ಅಂಗಿ} (aṅgi)}{ಅಂಗಿ (aṅgi)}]{\kn{ಅಂಗಿ} (aṅgi) — a shirt}

\label{lesson:KA-C120-angi}



\begin{quote}
Before the new one: say the Kannada for a ball, then the Kannada for a kite.
\end{quote}

\subsection*{You'll want to know: \kn{ಅಂಗಿ}}
\textbf{\kn{ಅಂಗಿ}} — \emph{aṅgi} — \textquotedblleft{}a shirt\textquotedblright{}.

\begin{grammarlens}[title={What you've built}]
One more word for this chapter.
\end{grammarlens}

\subsection*{Your turn}
\begin{itemize}
\raggedright
  \item \emph{Say it:} \emph{aṅgi}
  \item \emph{Say it:} \emph{aṅgi}, once more
  \item \emph{Say it:} say \emph{aṅgi}
  \item \emph{Recall:} say the Kannada for a ball, then the Kannada for a kite, then say \emph{aṅgi} again
\end{itemize}

\begin{tcolorbox}[breakable,colback=teal!4,colframe=teal!35!black,title={Before you move on}]
What does \kn{ಅಂಗಿ} mean? (\textquotedblleft{}A shirt\textquotedblright{}.) Say it once more.
\end{tcolorbox}
\section[\texorpdfstring{\kn{ಸೀರೆ} (sīre)}{ಸೀರೆ (sīre)}]{\kn{ಸೀರೆ} (sīre) — a sari}

\label{lesson:KA-C120-sire}



\begin{quote}
Before the new one: say the Kannada for a kite, then the Kannada for a shirt.
\end{quote}

\subsection*{You'll want to know: \kn{ಸೀರೆ}}
\textbf{\kn{ಸೀರೆ}} — \emph{sīre} — \textquotedblleft{}a sari\textquotedblright{}.

\begin{grammarlens}[title={What you've built}]
One more word for this chapter.
\end{grammarlens}

\subsection*{Your turn}
\begin{itemize}
\raggedright
  \item \emph{Say it:} \emph{sīre}
  \item \emph{Say it:} \emph{sīre}, once more
  \item \emph{Say it:} say \emph{sīre}
  \item \emph{Recall:} say the Kannada for a kite, then the Kannada for a shirt, then say \emph{sīre} again
\end{itemize}

\begin{tcolorbox}[breakable,colback=teal!4,colframe=teal!35!black,title={Before you move on}]
What does \kn{ಸೀರೆ} mean? (\textquotedblleft{}A sari\textquotedblright{}.) Say it once more.
\end{tcolorbox}
\section[\texorpdfstring{\kn{ಪಂಚೆ} (paṁce)}{ಪಂಚೆ (paṁce)}]{\kn{ಪಂಚೆ} (paṁce) — a dhoti}

\label{lesson:KA-C120-pamce}



\begin{quote}
Before the new one: say the Kannada for a shirt, then the Kannada for a sari.
\end{quote}

\subsection*{You'll want to know: \kn{ಪಂಚೆ}}
\textbf{\kn{ಪಂಚೆ}} — \emph{paṁce} — \textquotedblleft{}a dhoti\textquotedblright{}.

\begin{grammarlens}[title={What you've built}]
One more word for this chapter.
\end{grammarlens}

\subsection*{Your turn}
\begin{itemize}
\raggedright
  \item \emph{Say it:} \emph{paṁce}
  \item \emph{Say it:} \emph{paṁce}, once more
  \item \emph{Say it:} say \emph{paṁce}
  \item \emph{Recall:} say the Kannada for a shirt, then the Kannada for a sari, then say \emph{paṁce} again
\end{itemize}

\begin{tcolorbox}[breakable,colback=teal!4,colframe=teal!35!black,title={Before you move on}]
What does \kn{ಪಂಚೆ} mean? (\textquotedblleft{}A dhoti\textquotedblright{}.) Say it once more.
\end{tcolorbox}
\section[\texorpdfstring{\kn{ಚಪ್ಪಲಿ} (cappali)}{ಚಪ್ಪಲಿ (cappali)}]{\kn{ಚಪ್ಪಲಿ} (cappali) — sandals}

\label{lesson:KA-C120-cappali}



\begin{quote}
Before the new one: say the Kannada for a sari, then the Kannada for a dhoti.
\end{quote}

\subsection*{You'll want to know: \kn{ಚಪ್ಪಲಿ}}
\textbf{\kn{ಚಪ್ಪಲಿ}} — \emph{cappali} — \textquotedblleft{}sandals\textquotedblright{}.

\begin{grammarlens}[title={What you've built}]
One more word for this chapter.
\end{grammarlens}

\subsection*{Your turn}
\begin{itemize}
\raggedright
  \item \emph{Say it:} \emph{cappali}
  \item \emph{Say it:} \emph{cappali}, once more
  \item \emph{Say it:} say \emph{cappali}
  \item \emph{Recall:} say the Kannada for a sari, then the Kannada for a dhoti, then say \emph{cappali} again
\end{itemize}

\begin{tcolorbox}[breakable,colback=teal!4,colframe=teal!35!black,title={Before you move on}]
What does \kn{ಚಪ್ಪಲಿ} mean? (\textquotedblleft{}Sandals\textquotedblright{}.) Say it once more.
\end{tcolorbox}
